% Options for packages loaded elsewhere
\PassOptionsToPackage{unicode}{hyperref}
\PassOptionsToPackage{hyphens}{url}
\documentclass[
]{article}
\usepackage{xcolor}
\usepackage[margin=1in]{geometry}
\usepackage{amsmath,amssymb}
\setcounter{secnumdepth}{-\maxdimen} % remove section numbering
\usepackage{iftex}
\ifPDFTeX
  \usepackage[T1]{fontenc}
  \usepackage[utf8]{inputenc}
  \usepackage{textcomp} % provide euro and other symbols
\else % if luatex or xetex
  \usepackage{unicode-math} % this also loads fontspec
  \defaultfontfeatures{Scale=MatchLowercase}
  \defaultfontfeatures[\rmfamily]{Ligatures=TeX,Scale=1}
\fi
\usepackage{lmodern}
\ifPDFTeX\else
  % xetex/luatex font selection
\fi
% Use upquote if available, for straight quotes in verbatim environments
\IfFileExists{upquote.sty}{\usepackage{upquote}}{}
\IfFileExists{microtype.sty}{% use microtype if available
  \usepackage[]{microtype}
  \UseMicrotypeSet[protrusion]{basicmath} % disable protrusion for tt fonts
}{}
\makeatletter
\@ifundefined{KOMAClassName}{% if non-KOMA class
  \IfFileExists{parskip.sty}{%
    \usepackage{parskip}
  }{% else
    \setlength{\parindent}{0pt}
    \setlength{\parskip}{6pt plus 2pt minus 1pt}}
}{% if KOMA class
  \KOMAoptions{parskip=half}}
\makeatother
\usepackage{color}
\usepackage{fancyvrb}
\newcommand{\VerbBar}{|}
\newcommand{\VERB}{\Verb[commandchars=\\\{\}]}
\DefineVerbatimEnvironment{Highlighting}{Verbatim}{commandchars=\\\{\}}
% Add ',fontsize=\small' for more characters per line
\usepackage{framed}
\definecolor{shadecolor}{RGB}{248,248,248}
\newenvironment{Shaded}{\begin{snugshade}}{\end{snugshade}}
\newcommand{\AlertTok}[1]{\textcolor[rgb]{0.94,0.16,0.16}{#1}}
\newcommand{\AnnotationTok}[1]{\textcolor[rgb]{0.56,0.35,0.01}{\textbf{\textit{#1}}}}
\newcommand{\AttributeTok}[1]{\textcolor[rgb]{0.13,0.29,0.53}{#1}}
\newcommand{\BaseNTok}[1]{\textcolor[rgb]{0.00,0.00,0.81}{#1}}
\newcommand{\BuiltInTok}[1]{#1}
\newcommand{\CharTok}[1]{\textcolor[rgb]{0.31,0.60,0.02}{#1}}
\newcommand{\CommentTok}[1]{\textcolor[rgb]{0.56,0.35,0.01}{\textit{#1}}}
\newcommand{\CommentVarTok}[1]{\textcolor[rgb]{0.56,0.35,0.01}{\textbf{\textit{#1}}}}
\newcommand{\ConstantTok}[1]{\textcolor[rgb]{0.56,0.35,0.01}{#1}}
\newcommand{\ControlFlowTok}[1]{\textcolor[rgb]{0.13,0.29,0.53}{\textbf{#1}}}
\newcommand{\DataTypeTok}[1]{\textcolor[rgb]{0.13,0.29,0.53}{#1}}
\newcommand{\DecValTok}[1]{\textcolor[rgb]{0.00,0.00,0.81}{#1}}
\newcommand{\DocumentationTok}[1]{\textcolor[rgb]{0.56,0.35,0.01}{\textbf{\textit{#1}}}}
\newcommand{\ErrorTok}[1]{\textcolor[rgb]{0.64,0.00,0.00}{\textbf{#1}}}
\newcommand{\ExtensionTok}[1]{#1}
\newcommand{\FloatTok}[1]{\textcolor[rgb]{0.00,0.00,0.81}{#1}}
\newcommand{\FunctionTok}[1]{\textcolor[rgb]{0.13,0.29,0.53}{\textbf{#1}}}
\newcommand{\ImportTok}[1]{#1}
\newcommand{\InformationTok}[1]{\textcolor[rgb]{0.56,0.35,0.01}{\textbf{\textit{#1}}}}
\newcommand{\KeywordTok}[1]{\textcolor[rgb]{0.13,0.29,0.53}{\textbf{#1}}}
\newcommand{\NormalTok}[1]{#1}
\newcommand{\OperatorTok}[1]{\textcolor[rgb]{0.81,0.36,0.00}{\textbf{#1}}}
\newcommand{\OtherTok}[1]{\textcolor[rgb]{0.56,0.35,0.01}{#1}}
\newcommand{\PreprocessorTok}[1]{\textcolor[rgb]{0.56,0.35,0.01}{\textit{#1}}}
\newcommand{\RegionMarkerTok}[1]{#1}
\newcommand{\SpecialCharTok}[1]{\textcolor[rgb]{0.81,0.36,0.00}{\textbf{#1}}}
\newcommand{\SpecialStringTok}[1]{\textcolor[rgb]{0.31,0.60,0.02}{#1}}
\newcommand{\StringTok}[1]{\textcolor[rgb]{0.31,0.60,0.02}{#1}}
\newcommand{\VariableTok}[1]{\textcolor[rgb]{0.00,0.00,0.00}{#1}}
\newcommand{\VerbatimStringTok}[1]{\textcolor[rgb]{0.31,0.60,0.02}{#1}}
\newcommand{\WarningTok}[1]{\textcolor[rgb]{0.56,0.35,0.01}{\textbf{\textit{#1}}}}
\usepackage{graphicx}
\makeatletter
\newsavebox\pandoc@box
\newcommand*\pandocbounded[1]{% scales image to fit in text height/width
  \sbox\pandoc@box{#1}%
  \Gscale@div\@tempa{\textheight}{\dimexpr\ht\pandoc@box+\dp\pandoc@box\relax}%
  \Gscale@div\@tempb{\linewidth}{\wd\pandoc@box}%
  \ifdim\@tempb\p@<\@tempa\p@\let\@tempa\@tempb\fi% select the smaller of both
  \ifdim\@tempa\p@<\p@\scalebox{\@tempa}{\usebox\pandoc@box}%
  \else\usebox{\pandoc@box}%
  \fi%
}
% Set default figure placement to htbp
\def\fps@figure{htbp}
\makeatother
\setlength{\emergencystretch}{3em} % prevent overfull lines
\providecommand{\tightlist}{%
  \setlength{\itemsep}{0pt}\setlength{\parskip}{0pt}}
\usepackage{bookmark}
\IfFileExists{xurl.sty}{\usepackage{xurl}}{} % add URL line breaks if available
\urlstyle{same}
% fallback for those not using the hyperref driver hyperxmp:
\makeatletter
\@ifundefined{xmpquote}{\newcommand{\xmpquote}[1]{#1}}{}
\makeatother
\hypersetup{
  pdftitle={RExercise\#7},
  pdfauthor={Kent Jhon Herida},
  hidelinks,
  pdfcreator={LaTeX via pandoc}}

\title{RExercise\#7}
\author{Kent Jhon Herida}
\date{2026-11-06}

\begin{document}
\maketitle

\#\#\#Part A --- Creating Data Frames

\begin{Shaded}
\begin{Highlighting}[]
\NormalTok{respondents }\OtherTok{\textless{}{-}} \DecValTok{1}\SpecialCharTok{:}\DecValTok{20}
\NormalTok{sex }\OtherTok{\textless{}{-}} \FunctionTok{c}\NormalTok{(}\DecValTok{2}\NormalTok{, }\DecValTok{2}\NormalTok{, }\DecValTok{1}\NormalTok{, }\DecValTok{2}\NormalTok{, }\DecValTok{2}\NormalTok{, }\DecValTok{2}\NormalTok{, }\DecValTok{2}\NormalTok{, }\DecValTok{2}\NormalTok{, }\DecValTok{2}\NormalTok{, }\DecValTok{2}\NormalTok{, }\DecValTok{1}\NormalTok{, }\DecValTok{2}\NormalTok{, }\DecValTok{2}\NormalTok{, }\DecValTok{2}\NormalTok{, }\DecValTok{2}\NormalTok{, }\DecValTok{2}\NormalTok{, }\DecValTok{2}\NormalTok{, }\DecValTok{2}\NormalTok{, }\DecValTok{1}\NormalTok{, }\DecValTok{2}\NormalTok{)}
\NormalTok{fathers\_occ }\OtherTok{\textless{}{-}} \FunctionTok{c}\NormalTok{(}\DecValTok{1}\NormalTok{, }\DecValTok{3}\NormalTok{, }\DecValTok{3}\NormalTok{, }\DecValTok{3}\NormalTok{, }\DecValTok{1}\NormalTok{, }\DecValTok{2}\NormalTok{, }\DecValTok{3}\NormalTok{, }\DecValTok{1}\NormalTok{, }\DecValTok{1}\NormalTok{, }\DecValTok{1}\NormalTok{, }\DecValTok{3}\NormalTok{, }\DecValTok{2}\NormalTok{, }\DecValTok{1}\NormalTok{, }\DecValTok{3}\NormalTok{, }\DecValTok{3}\NormalTok{, }\DecValTok{1}\NormalTok{, }\DecValTok{3}\NormalTok{, }\DecValTok{1}\NormalTok{, }\DecValTok{2}\NormalTok{, }\DecValTok{1}\NormalTok{)}
\NormalTok{persons\_home }\OtherTok{\textless{}{-}} \FunctionTok{c}\NormalTok{(}\DecValTok{5}\NormalTok{, }\DecValTok{7}\NormalTok{, }\DecValTok{3}\NormalTok{, }\DecValTok{8}\NormalTok{, }\DecValTok{5}\NormalTok{, }\DecValTok{9}\NormalTok{, }\DecValTok{6}\NormalTok{, }\DecValTok{7}\NormalTok{, }\DecValTok{8}\NormalTok{, }\DecValTok{4}\NormalTok{, }\DecValTok{7}\NormalTok{, }\DecValTok{5}\NormalTok{, }\DecValTok{4}\NormalTok{, }\DecValTok{7}\NormalTok{, }\DecValTok{8}\NormalTok{, }\DecValTok{8}\NormalTok{, }\DecValTok{3}\NormalTok{, }\DecValTok{11}\NormalTok{, }\DecValTok{7}\NormalTok{, }\DecValTok{6}\NormalTok{)}
\NormalTok{siblings\_school }\OtherTok{\textless{}{-}} \FunctionTok{c}\NormalTok{(}\DecValTok{6}\NormalTok{, }\DecValTok{4}\NormalTok{, }\DecValTok{4}\NormalTok{, }\DecValTok{1}\NormalTok{, }\DecValTok{2}\NormalTok{, }\DecValTok{1}\NormalTok{, }\DecValTok{5}\NormalTok{, }\DecValTok{3}\NormalTok{, }\DecValTok{1}\NormalTok{, }\DecValTok{2}\NormalTok{, }\DecValTok{3}\NormalTok{, }\DecValTok{2}\NormalTok{, }\DecValTok{5}\NormalTok{, }\DecValTok{5}\NormalTok{, }\DecValTok{2}\NormalTok{, }\DecValTok{1}\NormalTok{, }\DecValTok{2}\NormalTok{, }\DecValTok{5}\NormalTok{, }\DecValTok{3}\NormalTok{, }\DecValTok{2}\NormalTok{)}
\NormalTok{house\_type }\OtherTok{\textless{}{-}} \FunctionTok{c}\NormalTok{(}\DecValTok{1}\NormalTok{, }\DecValTok{2}\NormalTok{, }\DecValTok{3}\NormalTok{, }\DecValTok{1}\NormalTok{, }\DecValTok{1}\NormalTok{, }\DecValTok{3}\NormalTok{, }\DecValTok{3}\NormalTok{, }\DecValTok{1}\NormalTok{, }\DecValTok{2}\NormalTok{, }\DecValTok{3}\NormalTok{, }\DecValTok{2}\NormalTok{, }\DecValTok{3}\NormalTok{, }\DecValTok{2}\NormalTok{, }\DecValTok{2}\NormalTok{, }\DecValTok{3}\NormalTok{, }\DecValTok{3}\NormalTok{, }\DecValTok{3}\NormalTok{, }\DecValTok{3}\NormalTok{, }\DecValTok{3}\NormalTok{, }\DecValTok{2}\NormalTok{)}

\NormalTok{respondents}
\end{Highlighting}
\end{Shaded}

\begin{verbatim}
##  [1]  1  2  3  4  5  6  7  8  9 10 11 12 13 14 15 16 17 18 19 20
\end{verbatim}

\begin{Shaded}
\begin{Highlighting}[]
\NormalTok{sex}
\end{Highlighting}
\end{Shaded}

\begin{verbatim}
##  [1] 2 2 1 2 2 2 2 2 2 2 1 2 2 2 2 2 2 2 1 2
\end{verbatim}

\begin{Shaded}
\begin{Highlighting}[]
\NormalTok{fathers\_occ}
\end{Highlighting}
\end{Shaded}

\begin{verbatim}
##  [1] 1 3 3 3 1 2 3 1 1 1 3 2 1 3 3 1 3 1 2 1
\end{verbatim}

\begin{Shaded}
\begin{Highlighting}[]
\NormalTok{persons\_home}
\end{Highlighting}
\end{Shaded}

\begin{verbatim}
##  [1]  5  7  3  8  5  9  6  7  8  4  7  5  4  7  8  8  3 11  7  6
\end{verbatim}

\begin{Shaded}
\begin{Highlighting}[]
\NormalTok{siblings\_school}
\end{Highlighting}
\end{Shaded}

\begin{verbatim}
##  [1] 6 4 4 1 2 1 5 3 1 2 3 2 5 5 2 1 2 5 3 2
\end{verbatim}

\begin{Shaded}
\begin{Highlighting}[]
\NormalTok{house\_type}
\end{Highlighting}
\end{Shaded}

\begin{verbatim}
##  [1] 1 2 3 1 1 3 3 1 2 3 2 3 2 2 3 3 3 3 3 2
\end{verbatim}

\begin{Shaded}
\begin{Highlighting}[]
\CommentTok{\# }\AlertTok{TODO}\CommentTok{: Before combining these into a data frame, confirm all six vectors}
\FunctionTok{length}\NormalTok{(respondents)}
\end{Highlighting}
\end{Shaded}

\begin{verbatim}
## [1] 20
\end{verbatim}

\begin{Shaded}
\begin{Highlighting}[]
\FunctionTok{length}\NormalTok{(sex)}
\end{Highlighting}
\end{Shaded}

\begin{verbatim}
## [1] 20
\end{verbatim}

\begin{Shaded}
\begin{Highlighting}[]
\FunctionTok{length}\NormalTok{(fathers\_occ)}
\end{Highlighting}
\end{Shaded}

\begin{verbatim}
## [1] 20
\end{verbatim}

\begin{Shaded}
\begin{Highlighting}[]
\FunctionTok{length}\NormalTok{(persons\_home)}
\end{Highlighting}
\end{Shaded}

\begin{verbatim}
## [1] 20
\end{verbatim}

\begin{Shaded}
\begin{Highlighting}[]
\FunctionTok{length}\NormalTok{(siblings\_school)}
\end{Highlighting}
\end{Shaded}

\begin{verbatim}
## [1] 20
\end{verbatim}

\begin{Shaded}
\begin{Highlighting}[]
\FunctionTok{length}\NormalTok{(house\_type)}
\end{Highlighting}
\end{Shaded}

\begin{verbatim}
## [1] 20
\end{verbatim}

\begin{Shaded}
\begin{Highlighting}[]
\CommentTok{\# have the same length using length() on each one (or wrap them in c()}
\CommentTok{\# and check for yourself that every value is 20)}

\FunctionTok{c}\NormalTok{(}
  \FunctionTok{length}\NormalTok{(respondents), }
  \FunctionTok{length}\NormalTok{(sex), }
  \FunctionTok{length}\NormalTok{(fathers\_occ), }
  \FunctionTok{length}\NormalTok{(persons\_home), }
  \FunctionTok{length}\NormalTok{(siblings\_school), }
  \FunctionTok{length}\NormalTok{(house\_type)}
\NormalTok{)}
\end{Highlighting}
\end{Shaded}

\begin{verbatim}
## [1] 20 20 20 20 20 20
\end{verbatim}

\begin{Shaded}
\begin{Highlighting}[]
\CommentTok{\# }\AlertTok{TODO}\CommentTok{: Combine all six vectors into a data frame called survey using}
\CommentTok{\# data.frame(), naming the columns: Respondent, Sex, FathersOccupation,}
\CommentTok{\# PersonsAtHome, SiblingsAtSchool, TypesOfHouses}
\CommentTok{\# Print the resulting data frame}
\NormalTok{survey }\OtherTok{\textless{}{-}} \FunctionTok{data.frame}\NormalTok{(}
  \AttributeTok{Respondent =}\NormalTok{ respondents,}
  \AttributeTok{Sex =}\NormalTok{ sex,}
  \AttributeTok{FathersOccupation =}\NormalTok{ fathers\_occ,}
  \AttributeTok{PersonsAtHome =}\NormalTok{ persons\_home,}
  \AttributeTok{SiblingsAtSchool =}\NormalTok{ siblings\_school,}
  \AttributeTok{TypesOfHouses =}\NormalTok{ house\_type}
\NormalTok{)}
\NormalTok{survey}
\end{Highlighting}
\end{Shaded}

\begin{verbatim}
##    Respondent Sex FathersOccupation PersonsAtHome SiblingsAtSchool
## 1           1   2                 1             5                6
## 2           2   2                 3             7                4
## 3           3   1                 3             3                4
## 4           4   2                 3             8                1
## 5           5   2                 1             5                2
## 6           6   2                 2             9                1
## 7           7   2                 3             6                5
## 8           8   2                 1             7                3
## 9           9   2                 1             8                1
## 10         10   2                 1             4                2
## 11         11   1                 3             7                3
## 12         12   2                 2             5                2
## 13         13   2                 1             4                5
## 14         14   2                 3             7                5
## 15         15   2                 3             8                2
## 16         16   2                 1             8                1
## 17         17   2                 3             3                2
## 18         18   2                 1            11                5
## 19         19   1                 2             7                3
## 20         20   2                 1             6                2
##    TypesOfHouses
## 1              1
## 2              2
## 3              3
## 4              1
## 5              1
## 6              3
## 7              3
## 8              1
## 9              2
## 10             3
## 11             2
## 12             3
## 13             2
## 14             2
## 15             3
## 16             3
## 17             3
## 18             3
## 19             3
## 20             2
\end{verbatim}

\begin{Shaded}
\begin{Highlighting}[]
\CommentTok{\#Describe the data. Get the structure or the summary of the data.}
\FunctionTok{str}\NormalTok{(survey)}
\end{Highlighting}
\end{Shaded}

\begin{verbatim}
## 'data.frame':    20 obs. of  6 variables:
##  $ Respondent       : int  1 2 3 4 5 6 7 8 9 10 ...
##  $ Sex              : num  2 2 1 2 2 2 2 2 2 2 ...
##  $ FathersOccupation: num  1 3 3 3 1 2 3 1 1 1 ...
##  $ PersonsAtHome    : num  5 7 3 8 5 9 6 7 8 4 ...
##  $ SiblingsAtSchool : num  6 4 4 1 2 1 5 3 1 2 ...
##  $ TypesOfHouses    : num  1 2 3 1 1 3 3 1 2 3 ...
\end{verbatim}

\begin{Shaded}
\begin{Highlighting}[]
\FunctionTok{summary}\NormalTok{(survey)}
\end{Highlighting}
\end{Shaded}

\begin{verbatim}
##    Respondent         Sex       FathersOccupation PersonsAtHome 
##  Min.   : 1.00   Min.   :1.00   Min.   :1.00      Min.   : 3.0  
##  1st Qu.: 5.75   1st Qu.:2.00   1st Qu.:1.00      1st Qu.: 5.0  
##  Median :10.50   Median :2.00   Median :2.00      Median : 7.0  
##  Mean   :10.50   Mean   :1.85   Mean   :1.95      Mean   : 6.4  
##  3rd Qu.:15.25   3rd Qu.:2.00   3rd Qu.:3.00      3rd Qu.: 8.0  
##  Max.   :20.00   Max.   :2.00   Max.   :3.00      Max.   :11.0  
##  SiblingsAtSchool TypesOfHouses
##  Min.   :1.00     Min.   :1.0  
##  1st Qu.:2.00     1st Qu.:2.0  
##  Median :2.50     Median :2.5  
##  Mean   :2.95     Mean   :2.3  
##  3rd Qu.:4.25     3rd Qu.:3.0  
##  Max.   :6.00     Max.   :3.0
\end{verbatim}

\begin{Shaded}
\begin{Highlighting}[]
\CommentTok{\#Check for the dimensions. How many rows and columns?}
\FunctionTok{dim}\NormalTok{(survey)}
\end{Highlighting}
\end{Shaded}

\begin{verbatim}
## [1] 20  6
\end{verbatim}

\begin{Shaded}
\begin{Highlighting}[]
\FunctionTok{nrow}\NormalTok{(survey)}
\end{Highlighting}
\end{Shaded}

\begin{verbatim}
## [1] 20
\end{verbatim}

\begin{Shaded}
\begin{Highlighting}[]
\FunctionTok{ncol}\NormalTok{(survey)}
\end{Highlighting}
\end{Shaded}

\begin{verbatim}
## [1] 6
\end{verbatim}

\begin{Shaded}
\begin{Highlighting}[]
\CommentTok{\#Create factors for Sex, Father\textquotesingle{}s Occupation and Types of Houses}
\NormalTok{survey }\OtherTok{\textless{}{-}} \FunctionTok{data.frame}\NormalTok{(}
  \AttributeTok{Respondent =}\NormalTok{ respondents,}
  \AttributeTok{Sex =}\NormalTok{ sex,}
  \AttributeTok{FathersOccupation =}\NormalTok{ fathers\_occ,}
  \AttributeTok{PersonsAtHome =}\NormalTok{ persons\_home,}
  \AttributeTok{SiblingsAtSchool =}\NormalTok{ siblings\_school,}
  \AttributeTok{TypesOfHouses =}\NormalTok{ house\_type}
\NormalTok{)}

\NormalTok{survey}\SpecialCharTok{$}\NormalTok{Sex }\OtherTok{\textless{}{-}} \FunctionTok{factor}\NormalTok{(}
\NormalTok{  survey}\SpecialCharTok{$}\NormalTok{Sex,}
  \AttributeTok{levels =} \FunctionTok{c}\NormalTok{(}\DecValTok{1}\NormalTok{, }\DecValTok{2}\NormalTok{),}
  \AttributeTok{labels =} \FunctionTok{c}\NormalTok{(}\StringTok{"Male"}\NormalTok{, }\StringTok{"Female"}\NormalTok{)}
\NormalTok{)}

\NormalTok{survey}\SpecialCharTok{$}\NormalTok{FathersOccupation }\OtherTok{\textless{}{-}} \FunctionTok{factor}\NormalTok{(}
\NormalTok{  survey}\SpecialCharTok{$}\NormalTok{FathersOccupation,}
  \AttributeTok{levels =} \FunctionTok{c}\NormalTok{(}\DecValTok{1}\NormalTok{, }\DecValTok{2}\NormalTok{, }\DecValTok{3}\NormalTok{),}
  \AttributeTok{labels =} \FunctionTok{c}\NormalTok{(}\StringTok{"Farmer"}\NormalTok{, }\StringTok{"Driver"}\NormalTok{, }\StringTok{"Others"}\NormalTok{)}
\NormalTok{)}

\NormalTok{survey}\SpecialCharTok{$}\NormalTok{TypesOfHouses }\OtherTok{\textless{}{-}} \FunctionTok{factor}\NormalTok{(}
\NormalTok{  survey}\SpecialCharTok{$}\NormalTok{TypesOfHouses,}
  \AttributeTok{levels =} \FunctionTok{c}\NormalTok{(}\DecValTok{1}\NormalTok{, }\DecValTok{2}\NormalTok{, }\DecValTok{3}\NormalTok{),}
  \AttributeTok{labels =} \FunctionTok{c}\NormalTok{(}\StringTok{"Wood"}\NormalTok{, }\StringTok{"Semi{-}Concrete"}\NormalTok{, }\StringTok{"Concrete"}\NormalTok{)}
\NormalTok{)}
\FunctionTok{str}\NormalTok{(survey)}
\end{Highlighting}
\end{Shaded}

\begin{verbatim}
## 'data.frame':    20 obs. of  6 variables:
##  $ Respondent       : int  1 2 3 4 5 6 7 8 9 10 ...
##  $ Sex              : Factor w/ 2 levels "Male","Female": 2 2 1 2 2 2 2 2 2 2 ...
##  $ FathersOccupation: Factor w/ 3 levels "Farmer","Driver",..: 1 3 3 3 1 2 3 1 1 1 ...
##  $ PersonsAtHome    : num  5 7 3 8 5 9 6 7 8 4 ...
##  $ SiblingsAtSchool : num  6 4 4 1 2 1 5 3 1 2 ...
##  $ TypesOfHouses    : Factor w/ 3 levels "Wood","Semi-Concrete",..: 1 2 3 1 1 3 3 1 2 3 ...
\end{verbatim}

\begin{Shaded}
\begin{Highlighting}[]
\NormalTok{survey}
\end{Highlighting}
\end{Shaded}

\begin{verbatim}
##    Respondent    Sex FathersOccupation PersonsAtHome SiblingsAtSchool
## 1           1 Female            Farmer             5                6
## 2           2 Female            Others             7                4
## 3           3   Male            Others             3                4
## 4           4 Female            Others             8                1
## 5           5 Female            Farmer             5                2
## 6           6 Female            Driver             9                1
## 7           7 Female            Others             6                5
## 8           8 Female            Farmer             7                3
## 9           9 Female            Farmer             8                1
## 10         10 Female            Farmer             4                2
## 11         11   Male            Others             7                3
## 12         12 Female            Driver             5                2
## 13         13 Female            Farmer             4                5
## 14         14 Female            Others             7                5
## 15         15 Female            Others             8                2
## 16         16 Female            Farmer             8                1
## 17         17 Female            Others             3                2
## 18         18 Female            Farmer            11                5
## 19         19   Male            Driver             7                3
## 20         20 Female            Farmer             6                2
##    TypesOfHouses
## 1           Wood
## 2  Semi-Concrete
## 3       Concrete
## 4           Wood
## 5           Wood
## 6       Concrete
## 7       Concrete
## 8           Wood
## 9  Semi-Concrete
## 10      Concrete
## 11 Semi-Concrete
## 12      Concrete
## 13 Semi-Concrete
## 14 Semi-Concrete
## 15      Concrete
## 16      Concrete
## 17      Concrete
## 18      Concrete
## 19      Concrete
## 20 Semi-Concrete
\end{verbatim}

\begin{Shaded}
\begin{Highlighting}[]
\CommentTok{\#Get the mean for all columns.}
\FunctionTok{colMeans}\NormalTok{(survey[}\FunctionTok{sapply}\NormalTok{(survey, is.numeric)])}
\end{Highlighting}
\end{Shaded}

\begin{verbatim}
##       Respondent    PersonsAtHome SiblingsAtSchool 
##            10.50             6.40             2.95
\end{verbatim}

\begin{Shaded}
\begin{Highlighting}[]
\CommentTok{\#Extract the 1st five rows and then all the columns using subsetting functions. What is the output?}
\NormalTok{survey[}\DecValTok{1}\SpecialCharTok{:}\DecValTok{5}\NormalTok{, ]}
\end{Highlighting}
\end{Shaded}

\begin{verbatim}
##   Respondent    Sex FathersOccupation PersonsAtHome SiblingsAtSchool
## 1          1 Female            Farmer             5                6
## 2          2 Female            Others             7                4
## 3          3   Male            Others             3                4
## 4          4 Female            Others             8                1
## 5          5 Female            Farmer             5                2
##   TypesOfHouses
## 1          Wood
## 2 Semi-Concrete
## 3      Concrete
## 4          Wood
## 5          Wood
\end{verbatim}

\begin{Shaded}
\begin{Highlighting}[]
\CommentTok{\#Select the variable types of houses then store the vector that results as types\_houses.}
\NormalTok{types\_houses }\OtherTok{\textless{}{-}}\NormalTok{ survey}\SpecialCharTok{$}\NormalTok{TypesOfHouses}
\NormalTok{types\_houses}
\end{Highlighting}
\end{Shaded}

\begin{verbatim}
##  [1] Wood          Semi-Concrete Concrete      Wood          Wood         
##  [6] Concrete      Concrete      Wood          Semi-Concrete Concrete     
## [11] Semi-Concrete Concrete      Semi-Concrete Semi-Concrete Concrete     
## [16] Concrete      Concrete      Concrete      Concrete      Semi-Concrete
## Levels: Wood Semi-Concrete Concrete
\end{verbatim}

\begin{Shaded}
\begin{Highlighting}[]
\CommentTok{\#Select only all Males respondent that their father occupation was farmer.}
\NormalTok{farmer\_males }\OtherTok{\textless{}{-}} \FunctionTok{subset}\NormalTok{(survey, Sex }\SpecialCharTok{==} \StringTok{"Male"} \SpecialCharTok{\&}\NormalTok{ FathersOccupation }\SpecialCharTok{==} \StringTok{"Farmer"}\NormalTok{)}
\NormalTok{farmer\_males}
\end{Highlighting}
\end{Shaded}

\begin{verbatim}
## [1] Respondent        Sex               FathersOccupation PersonsAtHome    
## [5] SiblingsAtSchool  TypesOfHouses    
## <0 rows> (or 0-length row.names)
\end{verbatim}

\begin{Shaded}
\begin{Highlighting}[]
\CommentTok{\#Select only all females respondent that have greater than or equal to 5 number of siblings attending school.}
\NormalTok{female\_many\_siblings }\OtherTok{\textless{}{-}} \FunctionTok{subset}\NormalTok{(survey, Sex }\SpecialCharTok{==} \StringTok{"Female"} \SpecialCharTok{\&}\NormalTok{ SiblingsAtSchool }\SpecialCharTok{\textgreater{}=} \DecValTok{5}\NormalTok{)}
\NormalTok{female\_many\_siblings}
\end{Highlighting}
\end{Shaded}

\begin{verbatim}
##    Respondent    Sex FathersOccupation PersonsAtHome SiblingsAtSchool
## 1           1 Female            Farmer             5                6
## 7           7 Female            Others             6                5
## 13         13 Female            Farmer             4                5
## 14         14 Female            Others             7                5
## 18         18 Female            Farmer            11                5
##    TypesOfHouses
## 1           Wood
## 7       Concrete
## 13 Semi-Concrete
## 14 Semi-Concrete
## 18      Concrete
\end{verbatim}

Question \#2: Before running data.frame(), you were asked to check that
all six vectors had the same length. Why is this check necessary
specifically for survey data like this, where every column represents an
answer from the SAME 20 respondents?

ans. Checking vector lengths ensures row-by-row data alignment across
respondents and prevents R from recycling shorter vectors or throwing
errors during data frame creation.

Question \#4:

In subset-logical, if you filter on Sex == 1 using the raw numeric code
rather than the word ``Male''.

Looking at the result, how could someone unfamiliar with the legend
misinterpret this output?

What problem does this illustrate about working with coded data before
it has been labeled?

(Filtering by raw numeric code 1 leaves the output showing only numbers,
which an unfamiliar reader could easily misinterpret as Female or
another category since the real-world meaning is missing, illustrating
how unlabeled coded data creates ambiguity and human error without clear
legend documentation.)

Question \#5:

Out of 20 respondents, how many are Male (Sex == 1)? Compute this using
nrow() on your filtered result from subset-logical or sum(survey\$Sex ==
1). Does this small sample size change how confident you would be
generalizing any pattern you find in this dataset to a larger
population? Explain your reasoning. (There are 3 Male respondents out of
the 20 total respondents.

Having only 3 males out of 20 total respondents significantly lowers
confidence in generalizing any pattern to a larger population because a
sample size that small lacks statistical power and introduces high
sampling variability. A tiny sample is easily distorted by individual
outliers and fails to represent the diversity of the broader population,
making any observed trends likely the result of random chance rather
than true societal patterns.)

Question \#6:

Run summary(). Every column shows numeric summary statistics (min, mean,
max, etc.) -- including Sex and TypesOfHouses, which are really
categories, not quantities.

Explain why R doesn't know, by itself, that these columns are
categorical, and what this tells you about the difference between how
data is stored versus what it means.

(R evaluates Sex and TypesOfHouses as numeric data because it only
tracks how the values are stored in memory, illustrating that computers
track raw storage types while analysts must explicitly define their
real-world categorical meaning.)

Question \#7:

Compare the mean of PersonsAtHome to the mean of Sex in your summary
output. Which of these two means is actually meaningful to report to
someone, and which one is not? Justify your answer.

(The mean of PersonsAtHome is meaningful because it measures a
quantitative count of household members, whereas the mean of Sex is
meaningless because numeric codes for categories are arbitrary labels
rather than measurable quantities.)

\end{document}
